\chapter{Conclusions and Future Work}
\label{chp:conclusions}

\section{Conclusions}
\label{sec:conclusions}

This thesis applied and extended established spurious feature detection and demographic
probing methodologies to Re-LAION-5B and Stable Diffusion v1.5, producing a systematic
empirical characterisation of how demographic bias is encoded, distributed, and amplified
across a web-scale training corpus and the generative model trained on it. The findings
address each of the five research objectives and collectively advance understanding of
bias mechanisms in large-scale vision systems.

\paragraph{Objective 1: Dataset acquisition, preparation, and annotation.}
A profession-conditioned subset of 200,000 images was retrieved from Re-LAION-5B across
200 occupations aligned with ISCO-08 categories, and a corresponding set of 200,000 images was generated using Stable Diffusion v1.5 with LCM‐LoRA acceleration for the standard generation pipeline. CLIP embeddings were recomputed from scratch using \texttt{openai/clip-vit-large-patch14} in the absence of pre-computed embeddings, and a two-stage structural validation pipeline was developed to enforce human-centric image suitability across both corpora. Demographic annotation was performed using the Realistic Gender Classifier and a VGG-16 MST10 Regression model selected through systematic evaluation against CCv2, MSTE, and FACET benchmarks, with skin tone extended to the full ten-point Monk Skin Tone scale beyond the binary and six-point scales used in prior work.

\paragraph{Objective 2: Spurious feature detection.}
The transformation-based framework of Zeng et al.\ was applied and extended to a
three-dataset configuration — Re-LAION-5B, Stable Diffusion v1.5, and COCO — enabling
per-class accuracy attribution that a binary design cannot provide. Fifteen
transformations spanning frequency, geometry, semantics, and colour were evaluated.
Every transformation exceeded the 33.3\% chance level, from unmodified images at 99.42\%
down to mean RGB at 64.19\%, establishing that dataset-specific spurious signatures
persist even under extreme information compression and are distributed across multiple
correlated visual attributes rather than concentrated in any single representation.

% \paragraph{Objective 3: Spurious feature correlations with demographics.}
% Demographic probing following the occlusion framework of Meister et al.\ revealed a
% fundamental asymmetry between the two datasets. In Re-LAION-5B, gender and skin tone
% information is localised entirely within the visible appearance of the depicted person:
% every occlusion variant collapses to chance, and no residual signal is recoverable from
% background context. In Stable Diffusion v1.5, demographic information is distributed
% across the full image, recoverable from background context alone (gender AUC 0.747; skin
% tone AUC 0.728), from person silhouette geometry (gender AUC 0.834), from object
% co-occurrence, skeletal geometry, and natural language descriptions (gender AUC 0.771).
% This cross-modal, cross-representational persistence demonstrates that the generative
% model does not merely inherit the demographic priors of its training corpus but projects
% them coherently onto the surrounding scene, spatial layout, and semantic content of
% generated images. Gender amplification was confirmed in 178 of 200 professions and skin
% tone amplification toward lighter tones in 175 of 200 professions, with the Dark-bin
% proportion collapsing by 83\% from Re-LAION-5B to Stable Diffusion v1.5.

\paragraph{Objective 3: Spurious feature correlations with demographics.}
Demographic probing following the occlusion framework of Meister et al.\ revealed a fundamental asymmetry between the two datasets. In Re-LAION-5B, gender and skin tone information is localised entirely within the visible appearance of the depicted person: every occlusion variant collapses to chance, and no residual signal is recoverable from background context. In Stable Diffusion v1.5, demographic information is uniquely distributed into scene context and background, recoverable from background context alone (gender AUC 0.747; skin tone AUC 0.728), from person silhouette geometry (gender AUC 0.834), from object co-occurrence, skeletal geometry, and natural language descriptions (gender AUC 0.771). This cross-modal, cross-representational persistence demonstrates that the generative model does not merely inherit the demographic priors of its training corpus but projects them coherently onto the surrounding scene, spatial layout, and semantic content of generated images. Gender amplification was confirmed in 178 of 200 professions and skin tone amplification toward lighter tones in 175 of 200 professions, with the Dark-bin proportion collapsing by 83\% from Re-LAION-5B to Stable Diffusion v1.5.

\paragraph{Objective 4: Inference-time debiasing.}
A combined pipeline of ITI-GEN demographic conditioning, ControlNet OpenPose structural
conditioning, and Chamfer latent-space colour alignment was implemented and evaluated on
a 10,000-image debiased doctor corpus. The pipeline successfully overrides the model's
strongly skewed learned prior, achieving a uniform gender and skin tone distribution by
construction. At the representation level, every non-person probe collapsed to at or near
chance for both gender and skin tone, including mean RGB (0.507 for gender; 0.482 for
skin tone), object detection, pose keypoints, depth at keypoints, and caption-based
probing (gender 0.530; skin tone 0.490). The near-chance mean RGB skin tone result
provides the most direct mechanistic evidence that the Chamfer colour alignment
successfully disrupts the photometric correlations that enable demographic signal to
propagate through colour statistics. Residual background-level demographic signal persists in the debiased corpus, with the background-retained probe (MaskRect) achieving AUC 0.995 for gender and 0.632 for skin tone, and the background-periphery probe (MaskRectNoBg) achieving AUC 0.701 for gender and 0.543 for skin tone, indicating that scene composition remains partially entangled with the demographic conditioning signal even under pose and appearance control.

\paragraph{Objective 5: Conclusions on bias mechanisms and mitigation.}
The combined evidence characterises a two-stage bias propagation process. Re-LAION-5B exhibits gender distributions that diverge from real-world labour force composition in a baseline-dependent manner, falling below US-specific norms in the majority of occupational groups while remaining broadly consistent with global workforce proportions in several others, and exhibits modest demographic signal distributed across colour, geometry, and semantic representations, but confines this signal to person-level appearance. Stable Diffusion v1.5 then amplifies
these imbalances substantially and transforms their distribution: the generative process
projects demographic priors beyond the person onto background, scene composition, object
layout, and textual description, producing a substantially more entrenched and
pervasive form of bias than is present in the training corpus. Inference-time debiasing
can measurably reduce spurious correlations at the representation level without model
retraining, but background-level entanglement persists and represents the primary
remaining challenge for this class of interventions.

\paragraph{Impact on the field.}
These findings have several direct implications for practitioners and researchers working
with large-scale visual datasets and generative systems. First, auditing a generative
model's demographic outputs is insufficient without also auditing the spurious feature
distributions of the training corpus, as the generative process amplifies and transforms
biases rather than simply reproducing them. Second, occlusion-based probing provides a
practical and interpretable diagnostic for distinguishing person-localised from
scene-distributed demographic signal, a distinction that aggregate distribution metrics
cannot detect. Third, the persistence of demographic signal in natural language
descriptions of generated images demonstrates that stereotyping in generative systems is
not confined to pixel statistics but extends into the semantic register, with implications
for any downstream application that consumes generated text-image pairs. Fourth, the
debiasing evaluation demonstrates that inference-time interventions can substantially reduce
low-level spurious correlations effectively, but background-level demographic
conditioning requires interventions that target scene composition directly rather than
person appearance alone.

\section{Future Work}
\label{sec:future_work}

Several directions present themselves as natural extensions of the work presented in this
thesis, spanning improvements to the debiasing methodology, broader evaluation scope, and
deeper investigation of the mechanisms identified here.

The most immediate extension concerns the debiasing pipeline. Residual background-level
demographic signal persists at AUC 0.995 for gender and 0.632 for skin tone under
background-only probing, indicating that scene composition remains partially entangled
with the demographic conditioning signal even under pose and appearance control. Future
work could investigate targeted background conditioning mechanisms that explicitly
decorrelate scene composition from demographic attributes, either through adversarial
background debiasing approaches or through extension of the ITI-GEN framework to
condition on background properties independently of person appearance. The weight editing
approach of Unified Concept Editing \cite{Gandikota_2024_UCE} offers a complementary
avenue, as its direct modification of cross-attention projections may suppress
background-level demographic associations that inference-time conditioning does not
address. A systematic comparison of conditioning-based approaches against weight editing
methods on the same probing protocol would establish whether background entanglement is
specific to the conditioning paradigm or is a general property of all inference-time
interventions.

The debiasing evaluation in this thesis is restricted to a single profession. Extending
the evaluation to a stratified sample spanning all ten ISCO-08 major groups would
determine whether the reduction of spurious correlations observed for the doctor
profession generalises across occupational contexts with different baseline demographic
distributions and different magnitudes of stereotyping. Professions at the extremes of
the observed skew distributions represent particularly informative test cases, as the
debiasing pipeline would be required to counteract the strongest learned priors in the
model.

The skin tone analysis in this thesis uses the full ten-point MST scale for the debiased
corpus, with labels derived from the generation configuration rather than from automated
classifiers. A complementary evaluation using validated automated skin tone estimation on
the debiased corpus would quantify the photometric fidelity of the conditioning,
assessing whether images conditioned on darker MST tokens exhibit measurably different
skin pigmentation statistics and whether the Chamfer colour alignment produces coherent
skin tone variation rather than global colour shifts. This would also enable the debiasing
evaluation to be compared directly with the eight-class standard corpus results.

The transformation-based probing framework applied in this work uses ConvNeXt-Tiny as
the classification backbone. Future work could evaluate whether the separability ceiling
identified here reflects genuine representational differences between datasets or the
capacity limitations of the chosen probe, by applying larger classification architectures
or self-supervised foundation model representations. The original Zeng et al.\ ablation
demonstrates limited impact of model size on dataset classification accuracy for
structural transformations, but this was conducted in a different dataset context and may
not generalise to the Re-LAION-5B versus Stable Diffusion v1.5 comparison.

Finally, the analysis is grounded in a subset of four Re-LAION-5B shards, covering
approximately 3.1\% of the full corpus. Extending the analysis to a larger shard sample
would determine whether the spurious feature profiles and demographic distributions
identified here are stable properties of the corpus as a whole or are influenced by
domain concentration effects specific to the retrieved subset. The link attrition and
domain concentration analyses conducted in this work establish that the retrieved subset
is non-uniform in its domain composition, and access to LAION crawl metadata would enable
a more principled characterisation of the representativeness of the surviving URL pool.